\section{增量式开发与环境管理}

\subsection{增量式开发的关键性}

在具备了 feature list 的基础上，Coding Agent 被要求每次只处理一个功能（one feature at a time）。这种增量式方法被证明是解决 Agent 贪多嚼不烂问题的关键。

增量式开发带来的好处包括：
\begin{itemize}
    \item 每个 session 的目标清晰、范围可控
    \item 即使某个 session 失败，损失也限制在单个功能范围内
    \item 通过 git commit 保留每个功能的实现记录，支持回滚
    \item 下一个 session 的 Agent 可以快速定位当前进展
\end{itemize}

\subsection{Session 结束时的清理工作}

每个 Coding Agent session 结束前必须执行以下清理操作：

\begin{enumerate}
    \item 将代码变更提交到 git，附带描述性的 commit message
    \item 在 \texttt{claude-progress.txt} 中写入本次 session 的工作摘要
    \item 更新 \texttt{feature\_list.json} 中已完成功能的状态
    \item 确保代码处于可编译、可运行的状态
\end{enumerate}

\begin{knowledgebox}{Git 作为安全网}
Git 不仅是版本控制工具，在这个系统中它还扮演着"安全网"的角色。当 Agent 的代码改动引入了 bug 或导致回归时，Agent 可以使用 \texttt{git revert} 回到之前的工作状态，而不是花大量时间调试。这种"先回滚再重来"的策略在长时间运行场景下尤为重要，因为每个 token 都是宝贵的资源。
\end{knowledgebox}

\subsection{Session 启动流程}

每个 Coding Agent session 启动时会执行一套标准化的流程来获取当前状态：

\begin{enumerate}
    \item 运行 \texttt{pwd} 确认工作目录
    \item 阅读 git log 和 progress 文件了解近期工作
    \item 阅读 feature list，选择下一个最高优先级的未完成功能
    \item 运行 \texttt{init.sh} 启动开发服务器
    \item 执行基本的端到端测试，确认现有功能正常工作
    \item 确认无回归 bug 后，开始实现新功能
\end{enumerate}

\begin{warningbox}{先测试再开发，而非先开发再测试}
一个关键的设计决策是：Agent 在开始新功能开发之前，必须先验证现有功能是否正常。如果跳过这一步直接开始新功能的开发，可能会在已有 bug 的基础上引入更多问题，使情况更加恶化。这与人类工程师的最佳实践一致：pull 最新代码后先跑一遍测试，确认 CI 是绿的再开始写新代码。
\end{warningbox}

\subsection{典型的 Session 启动消息}

以下是 Coding Agent 的典型启动行为序列：

\begin{lstlisting}[language={}, caption={Coding Agent 典型的 session 启动流程}]
[Assistant] I'll start by getting my bearings and understanding
            the current state of the project.
[Tool Use]  <bash - pwd>
[Tool Use]  <read - claude-progress.txt>
[Tool Use]  <read - feature_list.json>

[Assistant] Let me check the git log to see recent work.
[Tool Use]  <bash - git log --oneline -20>

[Assistant] Now let me check if there's an init.sh script
            to restart the servers.
<Starts the development server>

[Assistant] Let me verify that fundamental features are
            still working.
<Tests basic functionality>

[Assistant] Core features are working correctly. Let me
            review the feature list to determine what to
            implement next.
<Starts work on a new feature>
\end{lstlisting}

这种标准化的启动流程大大减少了每个 session 的"热身"时间，避免了 Agent 浪费 token 在"弄清楚怎么运行项目"这类重复性工作上。

\subsection{本章小结}

增量式开发的核心是每个 session 只处理一个功能，并在 session 结束时通过 git commit 和 progress 文件将环境留在干净状态。标准化的 session 启动流程确保每个 Agent 都能快速上手，先验证再开发的策略则防止了在已有问题上叠加新问题。


